% Generated draft: requires mathematical and source review.
\documentclass[11pt]{article}
\usepackage[margin=25mm]{geometry}
\usepackage{fontspec}
\setmainfont{Latin Modern Roman}
\usepackage{amsmath,amssymb,mathtools,unicode-math}
\setmathfont{Latin Modern Math}
\usepackage{enumitem,needspace,xurl,hyperref}
\setlength{\parindent}{0pt}
\setlength{\parskip}{5pt}
\setlength{\emergencystretch}{3em}
\newcommand{\N}{\mathbb{N}}
\newcommand{\Z}{\mathbb{Z}}
\newcommand{\Q}{\mathbb{Q}}
\newcommand{\R}{\mathbb{R}}
\newcommand{\C}{\mathbb{C}}
\newcommand{\F}{\mathbb{F}}
\newcommand{\A}{\mathbb{A}}
\newcommand{\PP}{\mathbb{P}}
\newcommand{\E}{\mathbb{E}}
\newcommand{\eps}{\varepsilon}
\newcommand{\dd}{\,\mathrm{d}}
\begin{document}
\section*{UnsolvedMath ID: 11000111}
\textbf{Source problem code:} catalogue entry
\subsection{Definitions and mathematical statement}
Let \(\Gamma\leq\operatorname{Mod}_g\) have finite index, in the source's surface convention. With trivial integral coefficients,
\[
 H^1(\Gamma;\Z)\cong\operatorname{Hom}(\Gamma,\Z).
\]
The cohomological reading is supported by the source's explicit phrase ``first cohomology'' and its final equivalent homomorphism question. Vanishing does not assert that the integral abelianization \(H_1(\Gamma;\Z)\) is zero; finite abelian quotients are not ruled out.

\emph{Problem 2.6.} Is \(H^1(\Gamma;\Z)=0\) for every finite-index subgroup of \(\operatorname{Mod}_g\)? The initial sentence does not state a genus restriction, but the source immediately gives counterexamples for \(g=1,2\) and focuses on the unresolved higher-genus setting. Those exceptions must remain explicit rather than extending the vanishing question indiscriminately to all genera.

The source reports these qualifications:
\begin{itemize}
\item For \(g\geq3\), McCarthy [37] proves vanishing if \(\Gamma\) contains the Torelli group.
\item For \(g=2\), McCarthy [37] and Taherkhani [43] give finite-index subgroups with nonzero first cohomology.
\item For each \(g\in\{1,2\}\) and integer \(n\geq2\), [31] gives a finite-index \(\Gamma_n\leq\operatorname{Mod}_g\) surjecting onto a rank-\(n\) free group. In particular, \(\operatorname{rank}H^1(\Gamma_n;\Z)\geq n\).
\item A positive answer to the unprovided Problem 1(a) would imply a positive answer here. The source says the \(g\geq3\) hypothesis in that Problem 1 is necessary, noting a finite-index subgroup of \(\operatorname{Mod}_2\) mapping onto a finite-index subgroup of \(\operatorname{Mod}_1\).
\end{itemize}
It asks separately the special case: for \(g\geq3\), if a finite-index \(\Gamma\) contains the Johnson kernel \(\mathcal K_g\), generated by Dehn twists about separating simple closed curves, must every homomorphism \(\Gamma\to\Z\) vanish? Containing \(\mathcal K_g\) is not replaced by the stronger requirement of containing the whole Torelli group.

The detached footnote records a later announcement by Behrstock and Margalit: an isomorphism between finite-index subgroups of \(\operatorname{Mod}_{1,2}^*\) need not be the restriction of an inner automorphism. Its citation is empty and the star convention is undefined; it is adjacent historical context, not a solution of the cohomology question.
\subsection{Short English statement}
Ask whether finite-index mapping-class subgroups in the higher-genus setting have no nonzero homomorphisms to \(\Z\), including the special case of subgroups containing the Johnson kernel. The source explicitly records genus-one/two counterexamples and a positive result when the full Torelli group is contained.
\subsection{Sources}
\begin{itemize}
\item Local UnsolvedMath registry, permanent ID 11000111, source code catalogue entry, Problem 2.6, PDF page 94. The trailing chapter header identifies M. Korkmaz; no individual proposer field is supplied.
\item Farb, \emph{Problems on Mapping Class Groups and Related Topics}. Supplied URL: \url{https://www.math.uchicago.edu/~farb/papers/mcgbook.pdf}; recorded access 2026-07-29; extraction pdf-font-structure. The volume contains 329 explicitly typeset headings across chapters, with repeated labels disambiguated by page.
\item McCarthy [37], Taherkhani [43], and the unnamed reference [31] are cited for the reported results. The local Problem 1(a) and its full statement are absent. The Behrstock--Margalit announcement in footnote 1 has an empty citation. No missing bibliographic entry is invented.
\item Supplied status open and OPEN-TRIAGE. The research summary says only ``Truncated.'' Current primary-literature status and redistribution terms for the publicly accessible author/publisher copy remain NEEDS\_REVIEW.
\end{itemize}
\end{document}
